% =========================================================
% Field and Order Axioms
% =========================================================

\subsection{Field and Order Axioms}
\label{sec:field-order-axioms}

\subsubsection{Definitions and Theorems}

The rational numbers are not merely a collection of equivalence classes.  Once addition and
multiplication are shown to be well-defined, $\mathbb{Q}$ becomes a field.  In addition, it carries a
natural order relation compatible with the field operations.  Together these two structures make
$\mathbb{Q}$ an \emph{ordered field}.

\paragraph{Field axioms.}
A field is a set $F$ with operations $+$ and $\cdot$ satisfying closure, associativity,
commutativity, additive and multiplicative identities, additive inverses, multiplicative inverses for
nonzero elements, and the distributive law.

\begin{remark}
The rational numbers satisfy these axioms under the usual operations inherited from their
construction.
\end{remark}

\paragraph{Order axioms.}
An ordered field is a field together with a relation $<$ satisfying trichotomy, transitivity,
compatibility with addition, and compatibility with multiplication by positive elements.

\begin{remark}
For many proofs it is more efficient to define order using the set of positive elements rather than to
take $<$ as primitive.
\end{remark}

\begin{definition}[Positive cone]
\label{def:positive-cone}
Let $F$ be a field.  A subset $P\subset F$ is called a \emph{positive cone} if it satisfies:
\begin{enumerate}
  \item if $a,b\in P$, then $a+b\in P$;
  \item if $a,b\in P$, then $ab\in P$;
  \item for every $a\in F$, exactly one of the following holds:
  \[
  a\in P,
  \qquad
  a=0,
  \qquad
  -a\in P.
  \]
\end{enumerate}
Define the order relation by
\[
a<b \iff b-a\in P.
\]
\end{definition}

\begin{remark}
This viewpoint makes the interaction between algebra and order especially transparent.  Statements
about inequalities can be translated into statements about positive elements and then handled
algebraically.
\end{remark}

\subsubsection{Immediate Consequences of the Positive Cone Axioms}

\begin{theorem}
\label{thm:squares-nonnegative}
For every $a\in F$,
\[
a^2\ge 0.
\]
\end{theorem}

\begin{proof}
By trichotomy, either $a\in P$, or $a=0$, or $-a\in P$.  In the first case $a^2\in P$ because
$P$ is closed under multiplication.  In the second case $a^2=0$.  In the third case
\[
(-a)(-a)=a^2\in P.
\]
Thus $a^2$ is positive or zero.
\end{proof}

\begin{theorem}
\label{thm:inverse-positive}
If $a>0$, then $a^{-1}>0$.
\end{theorem}

\begin{proof}
Suppose $a>0$.  If $a^{-1}<0$, then $-a^{-1}>0$.  Multiplying two positive elements gives
\[
a(-a^{-1})=-1>0.
\]
If also $1>0$, then trichotomy fails; so the assumption $a^{-1}<0$ is impossible.  Therefore
$a^{-1}>0$.
\end{proof}

\begin{theorem}
\label{thm:multiply-positive}
If $a<b$ and $c>0$, then
\[
ac<bc.
\]
\end{theorem}

\begin{proof}
From $a<b$ we have $b-a>0$.  Multiplying by $c>0$ gives
\[
c(b-a)>0.
\]
Thus
\[
bc-ac>0,
\]
which is equivalent to $ac<bc$.
\end{proof}

\begin{theorem}
\label{thm:multiply-negative}
If $a<b$ and $c<0$, then
\[
ac>bc.
\]
\end{theorem}

\begin{proof}
If $c<0$, then $-c>0$.  By the previous theorem,
\[
a(-c)<b(-c).
\]
This is
\[
-ac<-bc.
\]
Multiplying both sides by $-1$ reverses the inequality and yields $ac>bc$.
\end{proof}

\subsubsection{Positivity of $1$ and of the Natural Numbers}

\begin{theorem}
\label{thm:one-positive}
In an ordered field,
\[
1>0.
\]
\end{theorem}

\begin{proof}
By trichotomy exactly one of the following holds:
\[
1>0,
\qquad
1=0,
\qquad
-1>0.
\]
The second case is impossible in a field.  If $-1>0$, then multiplying two positive elements gives
\[
(-1)(-1)=1>0.
\]
Thus both $1$ and $-1$ would be positive, contradicting trichotomy.  Hence $1>0$.
\end{proof}

\begin{corollary}
\label{cor:n-positive}
For every natural number $n$,
\[
n>0.
\]
\end{corollary}

\begin{proof}
Each natural number is a finite sum of copies of $1$.  Since $1>0$ and the positive cone is closed
under addition, every such sum is positive.
\end{proof}

\begin{remark}[Archimedean property versus completeness]
\label{rem:arch-v-complete}
The positivity of $1$ anchors the order structure and leads naturally to the Archimedean property.
The Archimedean property holds in both $\mathbb{Q}$ and $\mathbb{R}$.  Completeness does not.  That is
why the rational numbers, although already an ordered field, still fail to capture all limits.
\end{remark}
